%% orchestrator-worker.tex: the orchestrator-worker pattern (Chapter 9).
%% House conceptual style; each worker is a `sub' node annotated with its own context.
%% _concept-style.tex: shared TikZ styles for the course's CONCEPTUAL block diagrams.
%% The leading underscore marks this as the one file in figures/ that is not a
%% figure; every other figures/*.tex holds exactly one figure body.
%%
%% \input this at the top of every conceptual figure (before \begin{tikzpicture})
%% so each such diagram speaks one visual language: rounded "block" nodes, stealth
%% arrows, and natural `<hue>!15' fills. Redefining these styles on each \input is
%% idempotent and harmless. The reference for the house parameters is
%% resources/STYLE-figures.md. Figure~\ref{figure:AgentLoop} is the exemplar.
%%
%% Scope: use for conceptual/relational diagrams (boxes-and-arrows). Do NOT use for
%% product mock-ups (the rig figures have their own shared language) or for
%% data/plot figures.
\tikzset{
  % the standard block node: geometry only — apply a `fill=<hue>!15' per node.
  conceptbox/.style={draw, rounded corners=4pt, minimum width=2.4cm,
                     minimum height=0.85cm, align=center, font=\small},
  % a flow arrow between blocks.
  conceptflow/.style={->, >=stealth', thick},
  % an edge/annotation label.
  conceptlabel/.style={font=\scriptsize, align=center},
}%

\begin{tikzpicture}[
    box/.style={conceptbox, minimum width=2.3cm, minimum height=1.0cm},
    sub/.style={conceptbox, minimum width=2.7cm, minimum height=0.9cm, font=\scriptsize},
    arr/.style={conceptflow}
]
\node[box, fill=orange!20] (orch) at (0,0) {Orchestrator};
\node[sub, fill=blue!12] (s1) at (4.7,1.7)  {sub-agent\\\emph{own context}};
\node[sub, fill=blue!12] (s2) at (4.7,0)    {sub-agent\\\emph{own context}};
\node[sub, fill=blue!12] (s3) at (4.7,-1.7) {sub-agent\\\emph{own context}};
\node[box, fill=green!20] (syn) at (9.4,0) {Synthesis};
\draw[arr] (orch) -- (s1);
\draw[arr] (orch) -- (s2) node[midway, above, font=\scriptsize]{task};
\draw[arr] (orch) -- (s3);
\draw[arr] (s1) -- (syn);
\draw[arr] (s2) -- (syn) node[midway, above, font=\scriptsize]{result};
\draw[arr] (s3) -- (syn);
\end{tikzpicture}%
